% Worked Examples: Concept 1.2.2 - Qualitative Analysis
\documentclass[11pt,a4paper]{article}
\usepackage{worked-examples-template}

\title{\textbf{Worked Examples}\\[0.5em]
\large Concept 1.2.2: Qualitative Analysis\\[0.3em]
\normalsize \coursesubtitle{} -- \coursetitle}
\author{Assoc.\ Prof.\ William Robertson \and Dr Sean McGowan}
\date{\today}

\begin{document}

\maketitle
\tableofcontents
\newpage

\section{Concept 1.2.2: Qualitative Analysis}

\subsection{Example 1: Phase Portrait of a Nonlinear Pendulum}

\paragraph{Problem:} Consider a simple pendulum with no damping:
\begin{align}
\ddot{\theta} + \omega_0^2 \sin\theta = 0
\end{align}

where $\omega_0 = \sqrt{g/\ell} = 2$ rad/s is the natural frequency, $g = 9.81$ m/s$^2$, and $\ell = 2.45$ m.

(a) Write the system as a set of first-order equations.

(b) Find all equilibrium points and classify them.

(c) Sketch the phase portrait showing trajectories in the $(\theta, \dot{\theta})$ plane.

(d) Identify which equilibria are stable and which are unstable.

\paragraph{Solution:}

\textbf{Step 1: First-order state-space form}

Define state variables:
\begin{align}
x_1 &= \theta \quad \text{(angle)} \\
x_2 &= \dot{\theta} \quad \text{(angular velocity)}
\end{align}

The system becomes:
\begin{align}
\dot{x}_1 &= x_2 \\
\dot{x}_2 &= -\omega_0^2 \sin x_1 = -4\sin x_1
\end{align}

\textbf{Step 2: Find equilibrium points}

At equilibrium, $\dot{x}_1 = \dot{x}_2 = 0$:
\begin{align}
x_2 &= 0 \\
-4\sin x_1 &= 0 \quad \Rightarrow \quad x_1 = n\pi, \quad n \in \mathbb{Z}
\end{align}

Equilibrium points:
\begin{itemize}
\item $(0, 0), (\pm 2\pi, 0), \ldots$ — pendulum hanging down (stable)
\item $(\pm \pi, 0), (\pm 3\pi, 0), \ldots$ — pendulum upright (unstable)
\end{itemize}

\textbf{Step 3: Linearize around equilibria}

The Jacobian matrix is:
\begin{align}
J = \begin{bmatrix} \frac{\partial \dot{x}_1}{\partial x_1} & \frac{\partial \dot{x}_1}{\partial x_2} \\ \frac{\partial \dot{x}_2}{\partial x_1} & \frac{\partial \dot{x}_2}{\partial x_2} \end{bmatrix} = \begin{bmatrix} 0 & 1 \\ -4\cos x_1 & 0 \end{bmatrix}
\end{align}

\textbf{At $(0, 0)$ (down position):}
\begin{align}
J|_{(0,0)} = \begin{bmatrix} 0 & 1 \\ -4 & 0 \end{bmatrix}
\end{align}

Eigenvalues: $\det(sI - J) = s^2 + 4 = 0 \Rightarrow s = \pm 2j$

Pure imaginary eigenvalues $\Rightarrow$ \textbf{center} (neutrally stable, periodic orbits)

\textbf{At $(\pi, 0)$ (up position):}
\begin{align}
J|_{(\pi,0)} = \begin{bmatrix} 0 & 1 \\ 4 & 0 \end{bmatrix}
\end{align}

Eigenvalues: $\det(sI - J) = s^2 - 4 = 0 \Rightarrow s = \pm 2$

Real eigenvalues with opposite signs $\Rightarrow$ \textbf{saddle point} (unstable)

\textbf{Step 4: Phase portrait description}

The phase portrait has the following features:

\begin{itemize}
\item \textbf{Centers at $\theta = 0, \pm 2\pi, \ldots$:} Closed elliptical trajectories around these points representing periodic oscillations (small angle: $\theta \approx \sin\theta$)
  
\item \textbf{Saddle points at $\theta = \pm\pi, \pm 3\pi, \ldots$:} Separatrix trajectories connecting saddle points, dividing the phase space into regions
  
\item \textbf{Closed orbits:} Trajectories near the stable equilibria represent oscillations with constant energy
  
\item \textbf{Open trajectories:} Trajectories passing through saddle points represent the pendulum rotating continuously (sufficient energy to go over the top)
  
\item \textbf{Energy contours:} Each trajectory corresponds to a constant total energy $E = \frac{1}{2}\dot{\theta}^2 - \omega_0^2\cos\theta$
\end{itemize}

\textbf{Behavior interpretation:}
\begin{itemize}
\item Small oscillations near $\theta = 0$: pendulum swings back and forth
\item Trajectories approaching saddle points: pendulum approaches upright position with decreasing velocity
\item Trajectories crossing saddle points horizontally: pendulum has exactly enough energy to reach the upright position
\item Wavy trajectories at large $|\dot{\theta}|$: pendulum rotates continuously, either clockwise or counterclockwise
\end{itemize}

\textbf{Key Insights:}
\begin{itemize}
\item Phase portraits provide qualitative understanding without solving equations analytically
\item Equilibrium points can be centers, saddles, nodes, or spirals
\item Nonlinear systems can have multiple equilibria with different stability properties
\item The separatrix divides fundamentally different behaviors (oscillation vs. rotation)
\item Energy is conserved in the undamped pendulum, leading to closed or periodic trajectories
\end{itemize}

\subsection{Example 2: Competing Species - Population Dynamics}

\paragraph{Problem:} Two species compete for the same resources. Their populations $x_1$ and $x_2$ evolve according to:
\begin{align}
\dot{x}_1 &= x_1(3 - x_1 - 2x_2) \\
\dot{x}_2 &= x_2(2 - x_1 - x_2)
\end{align}

(a) Find all equilibrium points.

(b) Linearize around each equilibrium and determine stability.

(c) Sketch the nullclines and phase portrait in the first quadrant ($x_1, x_2 \geq 0$).

\paragraph{Solution:}

\textbf{Step 1: Equilibrium points}

Set $\dot{x}_1 = \dot{x}_2 = 0$:
\begin{align}
x_1(3 - x_1 - 2x_2) &= 0 \\
x_2(2 - x_1 - x_2) &= 0
\end{align}

\textbf{Equilibria:}
\begin{itemize}
\item \textbf{$E_1 = (0, 0)$:} Both species extinct
\item \textbf{$E_2 = (3, 0)$:} Only species 1 survives (from $3 - x_1 = 0$)
\item \textbf{$E_3 = (0, 2)$:} Only species 2 survives (from $2 - x_2 = 0$)
\item \textbf{$E_4$:} Coexistence? Solve:
\begin{align}
3 - x_1 - 2x_2 &= 0 \\
2 - x_1 - x_2 &= 0
\end{align}
From the second equation: $x_1 = 2 - x_2$. Substitute into the first:
\begin{align}
3 - (2 - x_2) - 2x_2 = 0 \quad \Rightarrow \quad 1 + x_2 - 2x_2 = 0 \quad \Rightarrow \quad x_2 = 1
\end{align}

Then $x_1 = 2 - 1 = 1$. So $E_4 = (1, 1)$.
\end{itemize}

\textbf{Step 2: Linearization}

The Jacobian is:
\begin{align}
J = \begin{bmatrix} 3 - 2x_1 - 2x_2 & -2x_1 \\ -x_2 & 2 - x_1 - 2x_2 \end{bmatrix}
\end{align}

\textbf{At $E_1 = (0,0)$:}
\begin{align}
J|_{(0,0)} = \begin{bmatrix} 3 & 0 \\ 0 & 2 \end{bmatrix}
\end{align}
Eigenvalues: $\lambda_1 = 3$, $\lambda_2 = 2$ (both positive) $\Rightarrow$ \textbf{unstable node}

\textbf{At $E_2 = (3,0)$:}
\begin{align}
J|_{(3,0)} = \begin{bmatrix} -3 & -6 \\ 0 & -1 \end{bmatrix}
\end{align}
Eigenvalues: $\lambda_1 = -3$, $\lambda_2 = -1$ (both negative) $\Rightarrow$ \textbf{stable node}

\textbf{At $E_3 = (0,2)$:}
\begin{align}
J|_{(0,2)} = \begin{bmatrix} -1 & 0 \\ -2 & -2 \end{bmatrix}
\end{align}
Eigenvalues: $\lambda_1 = -1$, $\lambda_2 = -2$ (both negative) $\Rightarrow$ \textbf{stable node}

\textbf{At $E_4 = (1,1)$:}
\begin{align}
J|_{(1,1)} = \begin{bmatrix} -1 & -2 \\ -1 & -1 \end{bmatrix}
\end{align}
Eigenvalues: $\det(sI - J) = (s+1)^2 - 2 = s^2 + 2s - 1 = 0$
\begin{align}
s = \frac{-2 \pm \sqrt{4+4}}{2} = -1 \pm \sqrt{2}
\end{align}
$\lambda_1 = -1 + \sqrt{2} \approx 0.414 > 0$, $\lambda_2 = -1 - \sqrt{2} \approx -2.414 < 0$

One positive, one negative $\Rightarrow$ \textbf{saddle point} (unstable)

\textbf{Step 3: Nullclines and phase portrait}

\textbf{Nullclines} ($\dot{x}_1 = 0$ and $\dot{x}_2 = 0$):
\begin{itemize}
\item $\dot{x}_1 = 0$: $x_1 = 0$ or $x_1 + 2x_2 = 3$ (line from $(3,0)$ to $(0, 1.5)$)
\item $\dot{x}_2 = 0$: $x_2 = 0$ or $x_1 + x_2 = 2$ (line from $(2,0)$ to $(0, 2)$)
\end{itemize}

\textbf{Flow direction:}
\begin{itemize}
\item Below both nullclines: both populations grow
\item Above both: both decrease
\item Between: mixed behavior
\end{itemize}

\textbf{Outcome:} The system has two stable equilibria ($E_2$ and $E_3$) and one saddle ($E_4$). The separatrix through the saddle divides the phase space:
\begin{itemize}
\item Initial conditions on one side lead to $(3, 0)$ — species 1 wins
\item Initial conditions on the other side lead to $(0, 2)$ — species 2 wins
\item Coexistence $(1,1)$ is unstable
\end{itemize}

This is a classic example of \textbf{competitive exclusion}: both species cannot coexist in equilibrium.

\textbf{Key Insights:}
\begin{itemize}
\item Qualitative analysis reveals long-term behavior without explicit solutions
\item Nullclines divide the state space into regions with different flow directions
\item Competing species often lead to extinction of one species (competitive exclusion principle)
\item The basin of attraction determines which species wins based on initial populations
\end{itemize}
\end{document}
